\chapter{Deterministic OT Equal-Weighting Baseline}
\label{ch:bf-dpf-ot-baseline}

The previous chapter introduced the resampling bottleneck and explained why one
might replace discrete ancestor selection by a transport object.  The present
chapter slows that idea down.  Its purpose is to teach the deterministic optimal-
transport equal-weighting baseline itself before entropy, Sinkhorn iteration, or
learned acceleration are added.

The key conceptual move is simple.  A weighted particle cloud and an equal-weight
particle cloud are not the same object.  Deterministic OT equal-weighting asks for
a mass-redistribution rule that sends the weighted source cloud to an equal-weight
target cloud as cheaply as possible under a declared cost \citep{villani2003topics,peyre2019computational,reich2013nonparametric}.  This already changes
the resampling object relative to categorical ancestor sampling, but it does so in
a clean, finite-dimensional, deterministic way.

\section{What problem this chapter solves}
\label{sec:bf-ot-baseline-problem}

Suppose a particle filter has produced support points
\(x_1,\ldots,x_N\in\mathbb R^{d_x}\) with normalized weights
\(w_1,\ldots,w_N\).  A later stage of the filter may want an equally weighted
cloud so that no single particle dominates future propagation or summary
statistics.  Classical resampling achieves this by random duplication and
elimination \citep{doucet2001sequential,chopin2020introduction}.  The deterministic OT baseline instead asks:

\begin{quote}
How should mass be reassigned from the weighted cloud to an equal-weight output
cloud if we want to move it as little as possible under a declared transport
cost?
\end{quote}

This chapter answers that question in three layers:
\begin{enumerate}[label=(\roman*)]
  \item define the weighted and equal-weight empirical objects;
  \item define the coupling matrix that records how source mass is reassigned;
  \item convert that coupling into the equal-weight output cloud by barycentric
  projection.
\end{enumerate}
These are the same basic objects that underlie deterministic ensemble-transform
and OT-based equal-weighting formulations \citep{reich2013nonparametric,peyre2019computational}.

\begin{center}
\small
\begin{tabular}{p{0.20\textwidth} p{0.22\textwidth} p{0.22\textwidth} p{0.24\textwidth}}
\toprule
Stage & Source object & Output object & What later chapters change \\
\midrule
This chapter & weighted empirical cloud & deterministic coupling and barycentric equal-weight cloud & nothing yet; fixes the baseline object \\
Next chapter & same weighted cloud & entropically regularized coupling and barycentric cloud & replaces the deterministic LP by an entropic teacher and finite Sinkhorn route \\
Later OT chapters & same weighted cloud & accelerated, learned, or target-changing transports & may preserve the teacher, compress it, or replace it entirely \\
\bottomrule
\end{tabular}
\end{center}

\section{Running three-particle cell: deterministic OT version}
\label{sec:bf-ot-baseline-cell}

Reuse the running one-dimensional cell from the previous chapter:
\[
  x_1=-2,
  \qquad x_2=0,
  \qquad x_3=3,
  \qquad
  w=
  \left(
    \frac{1}{2},
    \frac{1}{3},
    \frac{1}{6}
  \right),
  \qquad
  u=
  \left(
    \frac{1}{3},
    \frac{1}{3},
    \frac{1}{3}
  \right).
\]
Using the source support itself as a simple square target support, the quadratic
cost matrix is
\begin{equation}
\label{eq:bf-ot-baseline-cost}
  C
  =
  \begin{pmatrix}
    0 & 4 & 25 \\
    4 & 0 & 9 \\
    25 & 9 & 0
  \end{pmatrix}.
\end{equation}
A feasible coupling for this problem is
\begin{equation}
\label{eq:bf-ot-baseline-feasible-coupling}
  \Pi^{\mathrm{cell}}
  =
  \begin{pmatrix}
    \frac{1}{3} & \frac{1}{6} & 0 \\
    0 & \frac{1}{6} & \frac{1}{6} \\
    0 & 0 & \frac{1}{6}
  \end{pmatrix}.
\end{equation}
The row sums are
\((1/2,1/3,1/6)\), matching the source weights, and the column sums are
\((1/3,1/3,1/3)\), matching the equal target marginal.  The barycentric output
cloud is therefore
\begin{align}
\label{eq:bf-ot-baseline-cell-barycentric}
  \tilde x^{(1)} &= 3\left(\frac{1}{3}(-2)\right) = -2, \\
  \tilde x^{(2)} &= 3\left(\frac{1}{6}(-2)+\frac{1}{6}(0)\right) = -1, \\
  \tilde x^{(3)} &= 3\left(\frac{1}{6}(0)+\frac{1}{6}(3)\right) = 1.5.
\end{align}
So one deterministic equal-weight cloud associated with this feasible transport is
\[
  \{-2,-1,1.5\}
  \quad\text{with equal weights }\frac{1}{3}.
\]
This single worked cell is enough to make three key points visible:
\begin{enumerate}[label=(\roman*)]
  \item the coupling is a mass-routing matrix, not yet the output cloud;
  \item the output particles need not be duplicates of the source support;
  \item deterministic equal-weighting already differs from categorical resampling
  even before entropy or differentiability are mentioned.
\end{enumerate}

\section{Objects and notation}
\label{sec:bf-ot-baseline-objects}

The chapter uses the following core objects.

\begin{center}
\small
\begin{tabular}{p{0.23\textwidth} p{0.64\textwidth}}
\toprule
Object & Meaning in this chapter \\
\midrule
$\hat\pi^N$ & weighted empirical source measure \\
$\hat\pi^N_{\mathrm{eq}}$ & equal-weight empirical output measure \\
$w$ & source weight vector \\
$u$ & target equal-weight vector, usually $(1/N,\ldots,1/N)$ \\
$C$ & declared transport cost matrix \\
$\Pi$ & coupling matrix sending source mass to target mass \\
$\mathcal U(w,u)$ & feasible coupling set with row sums $w$ and column sums $u$ \\
$\Pi_0^\star$ & optimal coupling for the declared unregularized OT problem \\
$B(\Pi)$ & barycentric equal-weight cloud induced by coupling $\Pi$ \\
\bottomrule
\end{tabular}
\end{center}

\section{Weighted cloud, equal-weight cloud, and the coupling set}
\label{sec:bf-ot-baseline-coupling}

The weighted source cloud is
\begin{equation}
\label{eq:bf-ot-baseline-weighted-measure}
  \hat\pi^N(dx)
  =
  \sum_{i=1}^N w_i\,\delta_{x_i}(dx),
  \qquad
  w_i\ge0,
  \quad
  \sum_{i=1}^N w_i=1.
\end{equation}
The equal-weight target cloud is
\begin{equation}
\label{eq:bf-ot-baseline-equal-measure}
  \hat\pi^N_{\mathrm{eq}}(dx)
  =
  \frac{1}{N}
  \sum_{j=1}^N \delta_{\tilde x^{(j)}}(dx).
\end{equation}
A coupling is a nonnegative matrix telling us how much mass from source particle
\(i\) contributes to output slot \(j\).  The feasible set is
\begin{equation}
\label{eq:bf-ot-baseline-feasible-set}
  \mathcal U(w,u)
  =
  \left\{
    \Pi\in\mathbb R_+^{N\times N}:
    \Pi\mathbf 1=w,
    \ \Pi'\mathbf 1=u
  \right\}.
\end{equation}
The row constraints say each source particle gives away exactly its original
weight.  The column constraints say each output slot receives exactly one equal
share of total mass.

\paragraph{Reader checkpoint.}
At this stage the coupling \(\Pi\) is only a bookkeeping object.  It is not yet a
random ancestor rule and not yet the final equal-weight cloud.  The next step is
to choose which feasible coupling is preferred.

\subsection{Why the filter needs a transport object, not just a scalar}
\label{sec:bf-ot-baseline-transport-object}

The scalable-OT survey sharpens the main practical boundary here.  For the
BayesFilter filtering recursion, a fast scalar optimal-transport value is not by
itself a usable resampling output.  The filter needs an object that actually acts
on the weighted particle cloud.  In the deterministic equal-weighting baseline,
that usable object is the coupling \(\Pi\) together with its barycentric action,
in the spirit of ensemble-transform and differentiable OT-resampling work
\citep{reich2013nonparametric,corenflos2021differentiable}.
Schematically, the required outputs are of the form
\begin{equation}
\label{eq:bf-ot-baseline-usable-object}
  \text{coupling matrix, barycentric projection, factorized transport operator,}
  \text{ or directly transported equal-weight cloud.}
\end{equation}
A reported scalar transport cost or value certificate may help compare candidate
transports, but it does not yet produce the next particle representation used by
the filter.

This is why the present chapter slows down on couplings and barycentric maps.
Given the weighted source measure
\begin{equation}
\label{eq:bf-ot-baseline-source-measure-repeat}
  \hat\pi^N(dx)
  =
  \sum_{i=1}^N w_i\,\delta_{x_i}(dx)
\end{equation}
and equal-weight target marginal \(u\), the mathematical output needed by later
code is not merely the scalar
\begin{equation}
\label{eq:bf-ot-baseline-scalar-only}
  \min_{\Pi\in\mathcal U(w,u)} \langle C,\Pi\rangle,
\end{equation}
but the minimizing coupling itself together with the map that turns that coupling
into an equal-weight cloud.  The later entropic and Sinkhorn chapter will replace
this exact linear-program teacher by an entropically regularized teacher, but it
inherits the same computational doctrine: a usable transport object must be
produced, not only a scalar score.

\section{The unregularized deterministic OT baseline}
\label{sec:bf-ot-baseline-unregularized}

Given a declared cost matrix \(C\), the deterministic transport baseline chooses a
feasible coupling with minimal cost:
\begin{equation}
\label{eq:bf-ot-baseline-lp}
  \Pi_0^\star
  \in
  \arg\min_{\Pi\in\mathcal U(w,u)}
  \langle C,\Pi\rangle.
\end{equation}
This is a linear program.  It is exact for the chosen deterministic OT object,
not for classical categorical resampling.  The distinction matters.  A reader who
keeps only one sentence in mind should keep this one:

\begin{quote}
The deterministic OT baseline solves a transport projection problem exactly for
that declared transport problem, but that problem is already different from
random ancestor sampling.
\end{quote}

\section{Barycentric projection: from coupling to equal-weight cloud}
\label{sec:bf-ot-baseline-barycentric}

Once a coupling has been chosen, the final equal-weight output cloud is obtained
by barycentric projection.  For each output slot \(j\), define
\begin{equation}
\label{eq:bf-ot-baseline-barycentric}
  B(\Pi)_j
  =
  \frac{1}{u_j}
  \sum_{i=1}^N \Pi_{ij}x_i.
\end{equation}
In the equal-weight case \(u_j=1/N\), this becomes
\begin{equation}
\label{eq:bf-ot-baseline-barycentric-equal}
  B(\Pi)_j
  =
  N\sum_{i=1}^N \Pi_{ij}x_i.
\end{equation}
Operationally, each column of the coupling is turned into an output particle by
forming the weighted average of the source particles that contribute mass to that
column.

This is the key step that makes the OT baseline implementation-facing.  The
coupling is a matrix of mass transfers; the barycentric map is the actual equal-
weight cloud later code would use.

\section{Algorithm: deterministic OT equal-weighting}
\label{sec:bf-ot-baseline-algorithm}

\begin{algorithm}[ht]
\caption{\texttt{deterministic\_ot\_equalweight\_step}}
\begin{algorithmic}[1]
\Require weighted particle cloud \(\{x_i,w_i\}_{i=1}^N\), equal target marginal \(u\), declared target support rule \(z_j\), and cost matrix \(C_{ij}=c(x_i,z_j)\)
\Ensure optimal coupling \(\Pi_0^\star\), barycentric equal-weight cloud \(\widetilde X\), and declared transport metadata
\State Form the feasible coupling set \(\mathcal U(w,u)\)
\State Solve the linear program \(\Pi_0^\star\in\arg\min_{\Pi\in\mathcal U(w,u)} \langle C,\Pi\rangle\)
\For{each output slot \(j=1,\ldots,N\)}
  \State Compute the barycentric output particle
  \[
    \tilde x^{(j)} = \frac{1}{u_j}\sum_{i=1}^N \Pi_{ij}^\star x_i
  \]
\EndFor
\State Return the equal-weight cloud \(\widetilde X=(\tilde x^{(1)},\ldots,\tilde x^{(N)})\), the coupling \(\Pi_0^\star\), the declared target support rule, and any transport diagnostics
\end{algorithmic}
\end{algorithm}

\paragraph{Approximation enters here.}
This algorithm is exact for the unregularized deterministic OT object in
\eqref{eq:bf-ot-baseline-lp}.  It is not exact for categorical resampling, and it
is not yet the entropically smoothed object used later for differentiable
computation.

\section{Filtering-loop insertion point}
\label{sec:bf-ot-baseline-loop}

Inside a filtering loop, this deterministic OT step sits after the weighted cloud
has been constructed and before the equal-weight cloud is passed to the next
propagation or summary stage.  In symbols:
\begin{equation}
\label{eq:bf-ot-baseline-loop}
  \{x_i,w_i\}_{i=1}^N
  \xrightarrow{\mathrm{deterministic\ OT\ equal\text{-}weighting}}
  \{\tilde x^{(j)},1/N\}_{j=1}^N.
\end{equation}
If auxiliary particle-local state is carried elsewhere in the filter, this chapter
does not yet decide how that state is transported.  That is a separate
implementation contract.  The present chapter fixes only the equal-weighting rule
for the particle locations and masses themselves.

\section{Verification and failure interpretation}
\label{sec:bf-ot-baseline-verification}

A minimal verification contract should check:
\begin{enumerate}[label=(\arabic*)]
  \item row sums equal the source weights exactly,
  \item column sums equal the equal target marginal exactly,
  \item the barycentric output has the declared shape \(N\times d_x\),
  \item the target support rule is recorded explicitly,
  \item any claim about downstream filtering is kept separate from this local OT
  equality check.
\end{enumerate}

A failed marginal check is evidence that the coupling solver or bookkeeping is
wrong.  It is not evidence that OT is conceptually invalid.  A good marginal check
is evidence only that the deterministic OT object has been implemented correctly,
not that the output cloud preserves the same law as categorical resampling.

\section{What this chapter does and does not claim}
\label{sec:bf-ot-baseline-boundary}

This chapter teaches the clean deterministic transport baseline for equal-
weighting.  It does \emph{not} claim:
\begin{itemize}
  \item that deterministic OT equal-weighting is the same object as classical
    categorical resampling,
  \item that it is automatically differentiable in the sense required later,
  \item that barycentric projection preserves the original particle-filter target,
  \item that the output cloud is already the final HMC or training target.
\end{itemize}
Its role is narrower: it fixes the exact deterministic OT baseline so that the
next chapter can explain what changes when entropy, Sinkhorn iteration, and
solver differentiation are added.  More precisely, the next chapter will preserve
the need for a usable transport object and the same barycentric output map while
changing the teacher coupling from the unregularized linear-program object to an
entropically regularized one.

\FloatBarrier
